\SourceLeaf{373}%
sen will, dass der Winkel \(MST\) nicht stumpf sein kann: allein hierbei ist doch ein wesentlicher Unterschied, weil in der That die Unmöglichkeit dieses Falles in aller Strenge bewiesen werden kann, welches weiter auszuführen aber hier nicht der Ort ist.

\Needspace{7\baselineskip}
\pdfbookmark[0]{Anzeigen. Müller: Theorie der Parallelen. 1822.}{anzeigen-mueller-1822}
\begin{center}
\rule{.85\textwidth}{.3pt}\par
Göttingische gelehrte Anzeigen.\quad 1822 October 28.\par
\rule{.85\textwidth}{.3pt}
\end{center}
\hspace*{\parindent}\textit{Theorie der Parallelen, von} \textsc{Carl Reinhard Müller}, \textit{Doctor der Philosophie, ausserordentlichem Professor der Mathematik u.\,s.\,w.} 40 S. in 4. Marburg 1822.

Rec. hat bereits vor sechs Jahren in diesen Blättern seine Überzeugung ausgesprochen, dass alle bisherigen Versuche, die Theorie der Parallellinien streng zu beweisen, oder die Lücke in der \textsc{Euklid}ischen Geometrie auszufüllen, uns diesem Ziele nicht näher gebracht haben, und kann nicht anders, als dies Urtheil auch auf alle späteren ihm bekannt gewordenen Versuche ausdehnen. Inzwischen bleiben doch manche solche Versuche, obgleich der eigentliche Hauptzweck verfehlt ist, wegen des darin bewiesenen Scharfsinns den Freunden der Geometrie lesenswerth, und. Rec. glaubt in dieser Rücksicht die vorliegende bei Gelegenheit einer Schulprüfung bekannt gemachte kleine Schrift besonders auszeichnen zu müssen. Den ganzen sinnreichen Ideengang des Verf. hier ausführlich darzulegen, wäre für unsre Blätter zu weitläufig und auch überflüssig, da die Schrift selbst gelesen zu werden verdient: aber sie hat ihre schwache Stelle, wie alle übrigen Versuche, und diese herauszuheben, ist der Zweck dieser Anzeige. Wir finden diese schwache Stelle S.\,15 in dem Beweise des Lehrsatzes des 15. Artikels. Dieser Lehrsatz ist der wahre Nerv der ganzen Theorie, welche fällt, sobald jener nicht streng bewiesen werden kann. Wir führen daher zuvörderst diesen Lehrsatz hier auf; die dazu gehörige Figur wird jeder leicht selbst zeichnen können.

Wenn jeder Winkel an der Grundlinie \(ON\) eines gleichschenkligen Dreiecks grösser ist, als der Winkel an der Spitze \(A\), und man setzt in \(O\) an die Seite \(OA\) einen Winkel von der Grösse des Winkels \(A\), dessen anderer Schenkel \(OL\) die \(AN\) in dem Punkte \(L\) zwischen \(A\) und \(N\) trifft, schneidet alsdann
